% Part: first-order-logic
% Chapter: introduction
% Section: semantic-notions

\documentclass[../../../include/open-logic-section]{subfiles}

\begin{document}

\olfileid{fol}{int}{sem}

\olsection{Semantische begrippen}

Hierboven merkten we op dat wanneer we nagaan of $\Sat{M}{!A}[s]$
geldt, we $s$ gemakshalve waarden aan alle !!{variable}s laten
toekennen, maar alleen de waarden gebruiken die het toekent aan de
!!{variable}s in~$!A$. In feite doen alleen de waarden van de
\emph{vrije} variabelen in~$!A$ ertoe. Omdat we zorgvuldig te werk
gaan, zullen we dit uiteraard bewijzen. !!^a{sentence} heeft geen vrije
variabelen; daarom maakt $s$ voor de vraag of zij in !!a{structure}
wordt vervuld helemaal niets uit. Als $!A$ !!a{sentence} is, kunnen we
$\Sat{M}{!A}$ dus definiëren als ``$\Sat{M}{!A}[s]$ voor alle~$s$''.
Dit blijkt waar te zijn dan en slechts dan als $\Sat{M}{!A}[s]$ voor
ten minste één~$s$ waar is. Voor een bruikbare definitie van vervulling
voor !!{formula}s moeten we variabelenbedelingen invoeren, maar voor
!!{sentence}s is vervulling onafhankelijk van die bedelingen.

Zodra ``$\Sat{M}{!A}$'' is gedefinieerd, weten we wat ``het geval'' en
``waar in'' betekenen voor !!{sentence}s van de eerste-ordelogica. Op
grond van de definitie van $\Sat{M}{!A}$ voor !!{sentence}s kunnen we
vervolgens de fundamentele semantische begrippen geldigheid, semantisch
gevolg en vervulbaarheid definiëren. !!^a{sentence} is geldig,
$\Entails !A$, als iedere !!{structure} haar vervult. Zij is een
semantisch gevolg van een verzameling !!{sentence}s,
$\Gamma \Entails !A$, als iedere !!{structure} die alle !!{sentence}s
in~$\Gamma$ vervult, ook~$!A$ vervult. Een verzameling !!{sentence}s is
vervulbaar als er !!a{structure} is die al haar !!{sentence}s
tegelijkertijd vervult.

Omdat !!{formula}s inductief zijn gedefinieerd en vervulling op haar
beurt door inductie naar de structuur van !!{formula}s is gedefinieerd,
kunnen we met inductie eigenschappen van onze semantiek bewijzen en de
gedefinieerde semantische begrippen met elkaar in verband brengen. We
zullen enkele van die eigenschappen formuleren en bewijzen. Dat doen we
deels omdat zij op zichzelf interessant zijn, maar vooral omdat veel
ervan van pas komen wanneer we het verband tussen semantiek en
!!{derivation}ssystemen onderzoeken. Daarvoor moeten we bovendien enkele
nog niet besproken syntactische begrippen en bewerkingen precies, dus
inductief, definiëren.

\end{document}
